\documentclass[UTF8,a4paper,12pt]{ctexart}

\usepackage[margin=2.45cm,headheight=15pt]{geometry}
\usepackage{booktabs,tabularx,array,longtable}
\usepackage{enumitem}
\usepackage{xcolor}
\usepackage{hyperref}
\usepackage{titlesec}
\usepackage{setspace}
\setCJKmainfont[Path=fonts/,AutoFakeBold=2]{NotoSansSC-VF.ttf}

\definecolor{ink}{HTML}{243447}
\definecolor{teal}{HTML}{287C78}
\definecolor{mist}{HTML}{EEF5F3}
\definecolor{gold}{HTML}{C98E38}
\hypersetup{colorlinks=true,linkcolor=teal,urlcolor=teal}
\setstretch{1.24}
\setlist[itemize]{leftmargin=2em,itemsep=3pt,topsep=4pt}
\setlist[enumerate]{leftmargin=2.2em,itemsep=4pt,topsep=4pt}
\titleformat{\section}{\Large\bfseries\color{ink}}{\thesection}{0.7em}{}
\titleformat{\subsection}{\large\bfseries\color{teal}}{\thesubsection}{0.7em}{}
\pagestyle{plain}

\begin{document}
\begin{titlepage}
  \centering
  \vspace*{2.8cm}
  {\large 项目开发计划书\par}
  \vspace{1.0cm}
  {\Huge\bfseries\color{ink} 灵感记录 APP\par}
  \vspace{4.8cm}
  \begin{minipage}[c][1.8cm][c]{0.82\textwidth}
    \centering\small
    让冒出来的想法，先有地方留下\\[5pt]
    少一点记录的阻力，多一点回看、串联与继续想下去的可能
  \end{minipage}
  \vfill
  \vspace*{1.3cm}
\end{titlepage}

\section*{项目摘要}
\addcontentsline{toc}{section}{项目摘要}
我们想做的不是又一个要求用户认真打字、分类、归档的备忘录。灵感常常出现在路上、聊天时，或者看到某个东西的一瞬间；等打开笔记、想好标题和目录，那个念头可能已经散了。这个项目希望把“先记下来”变得足够简单，再让 AI 帮用户把留下来的内容整理成能看、能找、能继续发展的灵感库。

首版围绕一条清楚的路径展开：用户用语音、文字或图片捕捉想法；系统保留原始内容，同时生成简短摘要、关键词和场景标签；用户回到“思维空间”浏览、搜索相关灵感，并按需要让 AI 做关联、提问或生成一个初步行动方案。AI 的整理结果可以修改，也可以撤回。它不能覆盖用户当时真正说过、拍过或写过的东西。

项目定位可以概括为“灵感记录 + AI 辅助整理”。我们希望守住两条产品原则：记录过程少阻力；AI 不急着把灵感修正成最标准、最安全、也最平庸的答案。首版以个人使用的闭环为中心，分享先控制在可选的单条导出，不把论坛、评论区或复杂社交放进核心范围。

\tableofcontents
\newpage

\section{项目概述}

\subsection{项目背景}
日常记录工具通常把整理工作留给用户：想标题、选分类、补上下文、日后再翻找。对待办事项和成熟笔记，这套方法有效；对一闪而过、还没成形的灵感，却容易让记录本身变成负担。

通用 AI 可以回答问题、生成内容，但如果在灵感刚出现时就把它交给模型，模型可能很快把独特的念头改写成常见方案。灵感有时需要先原样留下，允许它暂时不完整，之后再由本人决定要不要展开。

因此，本项目把重点放在“捕捉、整理、回访”三个环节：先留下原始想法，再提供可理解的组织方式，最后在合适的时候帮助用户重新看见它。

\subsection{项目目标}
\begin{enumerate}
  \item 做出一条可运行、可演示的个人灵感管理闭环：记录、整理、浏览、搜索和回访。
  \item 支持语音、文字和图片等低门槛输入；其中语音作为优先设计的入口，文字仍保留为补充和编辑方式。
  \item 让 AI 负责减轻整理负担，而不是替用户决定灵感的价值、方向或去留。
  \item 形成清楚、轻快、适合手机浏览的界面，并以可运行 Beta 程序和源代码作为阶段成果。
\end{enumerate}

\subsection{目标用户与使用场景}
首要用户是经常在学习、生活或创作过程中冒出零散念头，却不愿意为每个念头建立完整笔记的人。典型场景包括：走路时说下一句话；看到某个物件后拍照并补充想法；过几天通过关键词或场景标签找回旧记录；发现两条记录可能有关联时，决定是否把它们放到一起。

\subsection{产品原则}
\begin{itemize}
  \item \textbf{先收下，再整理：}新建记录尽量一步可达，允许内容不完整。
  \item \textbf{保留原样：}AI 摘要、合并和改写都不能替代原始录音、图片或文字。
  \item \textbf{建议由用户决定：}关联、归类、行动计划均提供确认、修改和忽略入口。
  \item \textbf{展示先于堆功能：}灵感库要便于浏览和回忆，不把所有材料挤成一长串文本。
  \item \textbf{核心优先：}先把个人记录体验做顺，再评估分享及其他扩展功能。
\end{itemize}

\section{产品方案与功能范围}

\subsection{核心使用流程}
\begin{enumerate}
  \item \textbf{捕捉：}打开应用后选择录音、文字或拍照；有条件时从桌面快捷入口直接开始记录。
  \item \textbf{整理：}保存后生成转写或摘要、关键词、时间等基础信息。AI 识别不确定时明确标出，不把猜测写成事实。
  \item \textbf{浏览：}在“思维空间”查看已整理的灵感卡片，可切换列表或简单关系视图；点开卡片可查看详情和原始证据。
  \item \textbf{继续：}用户可查看相关灵感、回答 AI 的启发式问题，或生成可编辑的下一步行动草案。
\end{enumerate}

\subsection{首版功能清单}
\begin{tabularx}{\textwidth}{>{\bfseries}p{2.6cm} X p{2.6cm}}
\toprule
功能 & 首版说明 & 优先级\\
\midrule
快速记录 & 录音、文字输入、图片附件；保存后立即反馈成功状态。 & 必做\\
原始证据 & 保留原始音频、文字和图片，记录时间；详情页可回看或播放。 & 必做\\
AI 整理 & 生成标题建议、简短摘要和关键词；生成内容可编辑、可重试。 & 必做\\
灵感空间 & 卡片/列表浏览、详情查看、基础筛选与搜索；呈现场景标签和材料类型。 & 必做\\
关联建议 & 根据关键词或语义提示可能相关的旧灵感，由用户确认后建立关系。 & 应做\\
发散与行动草案 & 提供少量引导问题，或把用户选中的灵感整理成可编辑步骤。 & 应做\\
桌面快捷入口 & 评估小组件或快捷方式的实现成本，选择至少一种快速启动方式。 & 争取\\
单条分享 & 将用户明确选中的灵感导出为图片或文本卡片；默认不公开。 & 可选\\
\bottomrule
\end{tabularx}

\subsection{明确暂缓的内容}
首版不建设论坛、评论区、关注关系和公共灵感广场，也不承诺后台整夜自动孵化。它们会带来内容治理、隐私和规模化问题，同时会分散团队对记录闭环和界面质量的投入。多端同步、手表端、复杂知识图谱和全自动合并也先列入后续评估。

\section{创新点与项目特色}

\begin{enumerate}
  \item \textbf{以“灵感”而非“文件”为组织对象。}系统围绕一条想法及其后续变化组织记录，尽量让用户看到的是能回忆起来的主题，而不是一堆未经整理的录音和附件。
  \item \textbf{整理结果与原始证据并存。}AI 可以概括和关联，但每条整理结果都能回到当时的录音、文字、图片和时间信息，避免“整理得很顺，却不知道原来怎么想的”。
  \item \textbf{AI 提供下一步，不抢走方向盘。}系统可提出问题、推荐关联或给出行动草案；用户可以接受、编辑、忽略，保留灵感暂时不成熟的空间。
  \item \textbf{把回看体验当作核心功能。}不同于只强调“记进去”，项目也关注如何用卡片、标签和简单关系视图，让用户愿意重新打开自己的灵感库。
\end{enumerate}

\section{技术方案}

\subsection{总体架构}
首版采用分层结构，客户端负责记录和浏览，服务端处理数据存储及 AI 调用，模型能力通过统一接口接入，避免产品逻辑绑定某一个模型供应方。

\begin{center}
\setlength{\tabcolsep}{4pt}
\begin{tabularx}{\textwidth}{>{\centering\arraybackslash}X c >{\centering\arraybackslash}X c >{\centering\arraybackslash}X}
\fbox{\parbox[c][1.6cm][c]{3.0cm}{\centering\small 客户端\\记录、浏览、确认}} & $\longrightarrow$ &
\fbox{\parbox[c][1.6cm][c]{3.0cm}{\centering\small 应用服务\\账号、灵感、权限}} & $\longrightarrow$ &
\fbox{\parbox[c][1.6cm][c]{3.0cm}{\centering\small AI 能力层\\转写、摘要、关联}} \\[8pt]
\multicolumn{5}{c}{\textcolor{teal}{\small 应用服务保存结构化记录与原始附件，并将处理结果返回客户端。}}
\end{tabularx}
\end{center}

\subsection{关键技术与选型原则}
\begin{itemize}
  \item \textbf{客户端：}会议讨论过 Flutter，也提到 MyGo 当前移动端适配尚未明确。建议第一周做一个小型验证：分别评估快速录音、图片选择、目标设备运行和界面开发效率，再由团队定案。若移动端能力是首版硬要求，优先选团队能稳定打包和调试的跨平台方案；不以尚未验证的适配能力作为进度前提。
  \item \textbf{服务端：}采用团队熟悉、便于接口联调的服务端方案。若项目已有 Go/MyGo 基础则沿用；否则先用成员能维护的简单服务，避免为技术栈本身增加学习和部署风险。
  \item \textbf{数据存储：}结构化保存灵感条目、标签、时间及关联关系；原始音频和图片作为独立附件保存，并以记录 ID 关联。首版不做复杂图数据库，关系视图由基础关联数据生成。
  \item \textbf{AI 接入：}将语音转写、摘要、标签和关联封装成独立能力接口。提示词和输出都要经过校验；失败时保留原始记录，允许稍后重试或手动编辑。
  \item \textbf{隐私与权限：}录音、相册和场景信息按需请求；定位、天气等环境信息不作为默认必需项。展示或分享前由用户明确选择内容。
\end{itemize}

\subsection{核心数据对象}
每条灵感至少包含：唯一标识、创建时间、输入类型、原始内容或附件地址、用户编辑内容、AI 摘要、关键词/标签、处理状态及用户确认的关联关系。AI 生成字段与用户原文分开保存；删除灵感时，应同时处理其附件和派生数据。

\section{可行性分析}

\subsection{技术可行性}
首版所需能力可拆为常见的移动端交互、文件上传与存储、文本检索和模型接口调用。主要技术风险不在于能否调用模型，而在于移动端录音及附件链路是否稳定、模型输出是否容易修改，以及团队能否在有限时间内完成端到端集成。因此先做最小链路验证，再扩展组织视图和快捷入口。

\subsection{进度与团队可行性}
会议已明确每位成员都参与代码，功能可由提出者或其他成员实现，并通过共享清单认领。团队可按客户端、服务端与数据、AI 能力、产品与测试等方向分工，但每项功能都需指定一位负责人，避免多人都以为别人会收尾。展示材料和影响团队理解的方案须由成员实际审阅后再对外使用。

\subsection{范围可行性}
在比赛周期内，完整实现多端、智能体、社交和自动孵化并不现实。若将范围收紧到“记录一条灵感—AI 整理—回看与搜索—确认关联”，可形成完整演示路径。桌面小组件、图谱展示和公开分享依次作为加分项，必要时可以从首版撤下。

\subsection{主要风险及应对}
\begin{tabularx}{\textwidth}{p{3.1cm} X X}
\toprule
风险 & 可能影响 & 应对方式\\
\midrule
移动端方案尚未定 & 录音、权限和打包工作延后 & 第一周做设备验证并冻结首版平台与框架；准备浏览器端演示备选。\\
模型结果不稳定 & 摘要失真或关联错误 & 保留原始证据；标注 AI 建议；所有合并和关联需用户确认。\\
功能范围持续增加 & 核心流程无法按时完成 & 用必做/应做/可选排序；新增功能经团队讨论并说明对排期的影响。\\
前端体验不足 & 产品看起来像功能拼盘 & 先制作关键页面原型；固定卡片层级、颜色和交互规则，再开发。\\
外部服务或网络不可用 & 演示中断、模型调用失败 & 提供本地样例数据和失败提示；准备不依赖实时模型的演示路径。\\
数据和权限顾虑 & 用户不愿意录音或上传材料 & 权限按需申请；给出清楚的删除入口；不默认公开个人内容。\\
\bottomrule
\end{tabularx}

\section{项目计划与排期}

以下按八周开发周期编排，具体日期应在比赛节点和团队课业安排确认后整体平移。每阶段结束都要有可检查的成果，不以“功能基本完成”作为验收标准。

\begin{longtable}{p{1.9cm} p{3.5cm} p{5.3cm} p{3.0cm}}
\toprule
阶段 & 重点工作 & 阶段产出 & 验收要点\\
\midrule
\endfirsthead
\toprule
阶段 & 重点工作 & 阶段产出 & 验收要点\\
\midrule
\endhead
\bottomrule
\endfoot
第 1 周 & 明确首版范围；验证客户端方案、录音与附件流程；确定接口和数据结构。 & 需求清单、页面草图、技术验证记录、任务认领表。 & 团队确认平台、必做功能和演示主线。\\
第 2 周 & 完成主要页面骨架、记录入口、灵感列表和详情页；搭建服务端基础接口。 & 可点击原型或早期运行版本；接口约定。 & 能打开应用并完成基本页面跳转。\\
第 3 周 & 打通文字记录、录音保存、图片附件和基础数据持久化。 & 可保存并重新打开的原始记录。 & 重启应用后记录仍可查看，附件不丢失。\\
第 4 周 & 接入语音转写和 AI 摘要、关键词；加入编辑与失败重试。 & 一条记录从输入到 AI 整理的完整链路。 & AI 失败不影响原始内容；结果可编辑。\\
第 5 周 & 完成灵感搜索、标签筛选和详情中的原始证据入口；实现关联建议。 & 可查找、可回溯、可确认关联的版本。 & 用户能从摘要返回原始记录，关系由用户确认。\\
第 6 周 & 打磨视觉与交互；加入引导式发散或行动草案；处理空状态和权限提示。 & 功能较完整的 Beta 候选版。 & 新用户能在少量提示下完成主流程。\\
第 7 周 & 集成测试、修复问题、准备样例数据；评估并补做一个低风险加分项。 & 稳定候选版、演示脚本、问题清单。 & 核心演示连续运行；失败场景有备用方案。\\
第 8 周 & 冻结功能、回归检查、整理源码和项目材料，录制展示视频。 & Beta 程序、源代码、计划书及展示材料。 & 材料内容与实际功能一致，团队完成人工审阅。\\
\bottomrule
\end{longtable}

\subsection{团队协作与变更}
采用共享功能清单认领任务，每项任务标明优先级、负责人、依赖和验收方式。提交代码前由本人检查生成内容及运行结果；涉及接口和数据结构的变更需提前告知相关成员。新增功能先讨论它解决什么问题、是否影响主流程和排期，再决定纳入或暂缓。每周至少一次集成检查，尽早发现模块之间的断点。

\section{验收标准与预期成果}

\subsection{首版验收标准}
\begin{itemize}
  \item 用户可以通过至少两种方式创建灵感，并在之后重新打开。
  \item 原始录音、文字或图片可以从详情页查看；AI 摘要不会覆盖原始材料。
  \item AI 能生成可编辑的摘要或关键词；失败时用户仍能保存和查看记录。
  \item 用户能用关键词或标签找到已有灵感，并对关联建议作出确认或忽略。
  \item 核心流程在目标设备上可连续演示；涉及权限、网络或模型失败时有明确提示。
\end{itemize}

\subsection{预期交付物}
\begin{enumerate}
  \item 可运行的 Beta 版本及完整源代码。
  \item 项目计划书、主要界面说明、技术结构和功能验收清单。
  \item 一段展示“快速记录—AI 整理—回看原始证据—发现关联”的演示视频。
  \item 测试记录与已知问题说明，标明尚未实现或仅为演示的功能。
\end{enumerate}

\section{结语}
这个项目不需要承诺“记下灵感就能改变人生”。更实际的目标，是让一条当下看来很小、甚至还不成熟的念头，不因为记录麻烦而消失；等用户想起它时，能找到当时的原始材料，也能决定要不要把它继续往前推。

首版是否成功，最终看两件事：记录是不是足够顺手，回头看时是不是比翻一堆录音和零散笔记更容易。先把这两件事做好，再谈更大的灵感网络。

\end{document}
